\section*{Conference Proceedings}
Graduate students and post-docs of my research group are indicated by * and $\dagger$, respectively.\\ 

\normalsize

\begin{etaremune}[topsep=0pt,itemsep=0pt,partopsep=0pt,parsep=0pt]

\item Roy Sarkar, D., Basu, S., Manuel, L., and Goswami, S. (2025). Uncertainty-aware optimization in engineered systems via gradient boosting and differential evolution, NASA and DNV UQ Challenge, in Proceedings of the 33rd European Safety and Reliability \& 33rd Society for Risk Analysis Europe Conference, \url{https://rpsonline.com.sg/proceedings/esrel-sra-e2025/html/toc.html}

\item Pierzyna, M., Saathof, R., and Basu, S. (2023). A multi-physics ensemble modeling framework for reliable $C_n^2$ estimation, in Proceedings of Environmental Effects on Light Propagation and Adaptive Systems VI, vol. 12731, SPIE, pp. 185–191. doi: 10.1117/12.2680997.

\item Basu, S., Watson, S. J., Cheneka, B. R., and Lacoa Arends, E (2020). Day-ahead wind power predictions at regional scales: Post-processing operational weather forecasts with a hybrid neural network, 17th International Conference on the European Energy Market (EEM20), IEEE, doi: 10.1109/EEM49802.2020.9221979.

\item Lacoa Arends, E., Watson, S. J., Basu, S., and Cheneka, B. R. (2020). Probabilistic wind power forecasting combining deep learning architectures, 17th International Conference on the European Energy Market (EEM20), IEEE, doi: 10.1109/EEM49802.2020.9221929.

\item Li, B., Basu, S., Watson, S. J., and Russchenberg, H. W. J. (2020). Quantifying the predictability of a `Dunkelflaute' event by utilizing a mesoscale model, Journal of Physics: Conference Series, 1618, 062042, doi: 10.1088/1742-6596/1618/6/062042.

\item Cheneka, B. R., Watson, S. J., and Basu, S. (2020). The impact of weather patterns on offshore wind power production, Journal of Physics: Conference Series, 1618, 062032, doi: 10.1088/1742-6596/1618/6/062032.

\item Lu, N.-Y., Basu, S., and Manuel, L (2017). Wind turbine loads during the evening transition period, 35th Wind Energy Symposium, AIAA SciTech Forum, (AIAA 2017-0681), doi: 10.2514/6.2017-0681.

\item Basu, S., and He$^{\dagger}$, P. (2015). Estimating refractive index structure parameter ($C_n^2$) profiles in the atmosphere: A wavelet transform-based approach, Proceedings of SPIE Optics \& Photonics,\\ doi: 10.1117/12.2188195.

\item He$^{\dagger}$, P. and Basu, S. (2015). Mesoscale modeling of optical turbulence ($C_n^2$) utilizing a novel physically-based parameterizations, Proceedings of SPIE Optics \& Photonics, doi: 10.1117/12.2188227.  

\item Wang*, Y., and Basu, S. (2014). Estimation of optical turbulence in the atmospheric surface layer from routine meteorological observations: An artificial neural network approach, Proceedings of SPIE Optics \& Photonics, doi: 10.1117/12.2063168. 

\item Nunalee*, C. G., He$^{\dagger}$, P., Basu, S., Vorontsov, M. A., and Fiorino, S. T. (2014). Impact of large-scale atmospheric refractive structures on optical wave propagation, Proceedings of SPIE Optics \& Photonics, doi: 10.1117/12.2063022. 

\item He*, P., Nunalee*, C. G., Basu, S., Vorontsov, M. A., and Fiorino, S. T. (2014). Current status and challenges in optical turbulence simulations in various layers of the Earth’s atmosphere, Proceedings of SPIE Optics \& Photonics, doi: 10.1117/12.2063023.

\item Basu, S., Nunalee*, C. G., He*, P., Fiorino, S. T., and Vorontsov, M. A. (2014). Reconstructing the prevailing meteorological and optical environment during the time of the Titanic disaster, Proceedings of SPIE Optics \& Photonics, doi: 10.1117/12.2063195. 

\item Nunalee*, C. G., and Basu, S. (2014). Mesoscale modeling of low-level jets over the North Sea, Wind Energy: Proceedings of the Euromech Colloquium, edited by Michael Hölling, Joachim Pienke and Stefan Ivanell, Springer, pages: 197-202. 

\item Basu, S., Holtslag, A. A. M., and Bosveld, F. C. (2011). GABLS3-LES intercomparison study, Workshop Proceedings – ECMWF Workshop on Diurnal Cycles and the Stable Boundary Layer, 7-10 November, Reading, UK, pages 75-82. 

\item Holtslag, A. A. M., Svensson, G., Basu, S., Beare, R. J., Bosveld, F. C., and Cuxart, J. (2011). Overview of the GEWEX Atmospheric Boundary Layer Study (GABLS), Workshop Proceedings – ECMWF Workshop on Diurnal Cycles and the Stable Boundary Layer, 7-10 November, Reading, UK, pages 11-24.

\item Moene, A. F., Baas, P., Bosveld, F. and Basu, S. (2011). LES model intercomparisons for the stable atmospheric boundary layer, Quality and Reliability of Large-Eddy Simulations II, ERCOFTAC Series, Vol. 16, Salvetti, M. V., Geurts, B., Meyers, J., Sagaut, P. (Eds.), Springer.  

\end{etaremune}
